% =====================================================================
%  Banco complementar --- Conceitos Iniciais e Materiais Magneticos
%  REVISADO. Substitui a versao gerada por template (10 clones em N2,
%  9 em N3 e 8 questoes conceituais de uma linha marcadas como N4).
%  O cap11 ja cobre: inseparabilidade dos polos (MC), definicao de
%  ferromagnetico por mu_r (MC), comparacao de dois materiais por B sob
%  mesmo H, classificacao para/diamagnetica a partir de mu_r, por que
%  ferromagneticos sao usados em nucleos, comparacao ar x ferro-silicio
%  no mesmo enrolamento e a descricao qualitativa do ciclo de histerese.
%  A LAMINACAO ja foi tratada quantitativamente no banco de Lenz, e o
%  ENTREFERRO no banco de fluxo -- por isso as N4 antigas sobre esses
%  dois temas foram removidas daqui.
%  Este banco explora: H, B e M como grandezas distintas, curva de
%  magnetizacao inicial e mu diferencial, temperatura de Curie, energia
%  do ciclo, materiais moles x duros com numeros, ponto de operacao em
%  nucleo nao linear, circuito magnetico com dois materiais e ensaio
%  experimental do ciclo.
%  Distribuicao mantida: 4 N1 + 11 N2 + 10 N3 + 8 N4 = 33
% =====================================================================

\subsubsection*{Banco complementar --- Conceitos iniciais e materiais magnéticos}

\begin{atencao}[Organização do banco]
Três grandezas convivem e são frequentemente confundidas. $\vec H$
($\si{\ampere\per\meter}$) é o campo \emph{aplicado}, determinado apenas pelas
correntes de condução; $\vec M$ é a \emph{magnetização} do material; e
$\vec B=\mu_0(\vec H+\vec M)$ é o que efetivamente age sobre cargas e produz
fluxo. Em materiais lineares, $\vec B=\mu_r\mu_0\vec H$ --- mas ferromagnéticos
\textbf{não} são lineares, e neles $\mu_r$ é uma função, não uma constante. O N4
trata de energia do ciclo, ponto de operação e ensaio.
\end{atencao}

\blocofixacao

\begin{questao}[1]{}
Um material linear tem $\mu_r=500$ e é submetido a
$H=\SI{100}{\ampere\per\meter}$. Determine $B$.
\dados{$\mu_0=4\pi\times10^{-7}\,\si{\tesla\meter\per\ampere}$}
\resposta{$B=\SI{62.8}{\milli\tesla}$}
\resolucao{
\[
B=\mu_r\mu_0H=(500)(4\pi\times10^{-7})(100)=\pot{6.28}{-2}\,\si{\tesla}.
\]
\textbf{Sem o material} ($\mu_r=1$), o mesmo $H$ produziria apenas
$\SI{126}{\micro\tesla}$. O material multiplicou $B$ por $500$ --- e é
exatamente esse o papel de um núcleo magnético.\\
\textbf{Atenção:} $H$ é fixado pelas correntes; $B$ é o que o material
\emph{responde}. Trocar o núcleo muda $B$, mas não muda $H$.}
\end{questao}

\begin{questao}[1]{}
Determine a magnetização $M$ do material da questão anterior.
\resposta{$M=\SI{49.9}{\kilo\ampere\per\meter}$}
\resolucao{De $B=\mu_0(H+M)$:
\[
M=\frac{B}{\mu_0}-H=\mu_rH-H=(\mu_r-1)H
=(499)(100)=\pot{4.99}{4}\,\si{\ampere\per\meter}.
\]
\textbf{Interpretação:} a magnetização do material equivale a uma corrente
superficial de $\SI{49.9}{\kilo\ampere\per\meter}$ --- quinhentas vezes maior
que a corrente de condução que a provocou.\\
\textbf{É por isso que o núcleo "amplifica":} ele não cria energia, apenas
alinha os momentos magnéticos já existentes no material, cujo efeito somado
supera em muito o do enrolamento.}
\end{questao}

\begin{questao}[1]{}
Define-se a suscetibilidade magnética como $\chi=\mu_r-1$. Classifique os
materiais com $\chi=+\pot{5}{-4}$ e $\chi=-\pot{2}{-5}$.
\resposta{Paramagnético e diamagnético, respectivamente}
\resolucao{\textbf{Paramagnético} ($\chi>0$, pequeno): os momentos atômicos
alinham-se \emph{fracamente} com o campo, reforçando-o de forma quase
imperceptível. Exemplos: alumínio, platina, oxigênio líquido.\\
\textbf{Diamagnético} ($\chi<0$): o campo induz momentos \emph{opostos}, e o
material é levemente repelido. Exemplos: cobre, água, bismuto, grafite. Todo
material tem componente diamagnética; ela só é observável quando não há
efeito maior a mascará-la.\\
\textbf{Ordens de grandeza:} $|\chi|\sim10^{-5}$ a $10^{-4}$ nesses dois casos,
contra $\chi\sim10^{3}$ em ferromagnéticos --- diferença de sete ordens de
grandeza.}
\end{questao}

\begin{questao}[1]{}
Um enrolamento de $500$ espiras percorrido por $\SI{0.4}{\ampere}$ envolve um
núcleo de caminho médio $\SI{25}{\centi\meter}$. Determine $H$.
\resposta{$H=\SI{800}{\ampere\per\meter}$}
\resolucao{Da lei de Ampère aplicada ao caminho médio,
\[
H=\frac{NI}{\ell}=\frac{500(0{,}4)}{0{,}25}=\SI{800}{\ampere\per\meter}.
\]
\textbf{Observação central:} $H$ \textbf{não depende do material}. Ele é
determinado exclusivamente pelas correntes e pela geometria.\\
Trocar o núcleo de ar por ferro não altera $H$ --- altera $B$, e portanto o
fluxo. É essa separação que torna útil distinguir as duas grandezas.}
\end{questao}

\blocoaplicacao

\begin{questao}[2]{}
Um núcleo de $\SI{10}{\centi\meter}$ de caminho médio, seção
$\SI{1}{\centi\meter\squared}$ e $\mu_r=1000$ recebe $\SI{200}{}$
ampère-espiras.
\begin{itensq}
\item Determine a relutância.
\item Determine o fluxo e a densidade de fluxo.
\item Avalie o resultado.
\end{itensq}
\resposta{$\mathcal{R}=\pot{7.96}{5}\,\si{\per\henry}$;
$\Phi=\SI{0.251}{\milli\weber}$; $B=\SI{2.51}{\tesla}$ --- \textbf{saturado}}
\resolucao{(a)
\[
\mathcal{R}=\frac{\ell}{\mu_r\mu_0A}
=\frac{0{,}10}{(1000)(4\pi\times10^{-7})(\pot{1}{-4})}
=\pot{7.96}{5}\,\si{\per\henry}.
\]
(b)
\[
\Phi=\frac{NI}{\mathcal{R}}=\frac{200}{\pot{7.96}{5}}
=\pot{2.51}{-4}\,\si{\weber};
\qquad
B=\frac{\Phi}{A}=\SI{2.51}{\tesla}.
\]
(c) \textbf{Resultado impossível.} Nenhum aço comum sustenta
$\SI{2.51}{\tesla}$ --- todos saturam entre $1{,}5$ e $\SI{2}{\tesla}$.\\
\textbf{O que realmente ocorreria:} conforme $B$ se aproxima da saturação,
$\mu_r$ desaba de $1000$ para valores próximos de $1$. A relutância dispara e o
fluxo estaciona em torno de $B_{sat}A=\SI{0.15}{\milli\weber}$.\\
\textbf{Regra de método:} calcular com $\mu_r$ constante e \emph{depois}
verificar $B$ contra $B_{sat}$. Se o valor exceder, o modelo linear não se
aplica e é preciso recorrer à curva $B\times H$ real.}
\end{questao}

\begin{questao}[2]{}
Um mesmo enrolamento produz $H=\SI{500}{\ampere\per\meter}$. Determine $B$ para
núcleo de ar, de alumínio ($\mu_r=1{,}00002$) e de ferro-silício
($\mu_r=4000$).
\resposta{$\SI{628}{\micro\tesla}$; $\SI{628}{\micro\tesla}$;
$\SI{2.51}{\tesla}$}
\resolucao{
\[
\text{ar}:\ B=\mu_0(500)=\SI{628}{\micro\tesla};
\]
\[
\text{alumínio}:\ B=(1{,}00002)\mu_0(500)=\SI{628}{\micro\tesla};
\]
\[
\text{ferro-silício}:\ B=(4000)\mu_0(500)=\SI{2.51}{\tesla}
\ \text{(saturaria antes)}.
\]
\textbf{Duas leituras.}\\
\emph{(i)} O alumínio é \emph{indistinguível} do ar para fins práticos --- sua
suscetibilidade de $\pot{2}{-5}$ altera $B$ na quinta casa decimal. É por isso
que se pode enrolar bobinas em carretéis de alumínio sem consequência magnética.\\
\emph{(ii)} O ferro-silício produziria $\SI{2.51}{\tesla}$ \emph{se fosse
linear} --- mas ele satura antes. Com $H=\SI{500}{\ampere\per\meter}$, um
ferro-silício real já está bem dentro do joelho da curva, com $B$ real em torno
de $1{,}5$ a $\SI{1.7}{\tesla}$.}
\end{questao}

\begin{questao}[2]{}
Um material apresenta $B=\SI{1.2}{\tesla}$ para $H=\SI{400}{\ampere\per\meter}$
e $B=\SI{1.5}{\tesla}$ para $H=\SI{1200}{\ampere\per\meter}$.
\begin{itensq}
\item Determine $\mu_r$ em cada ponto.
\item Interprete a variação.
\end{itensq}
\resposta{$\mu_r=2390$ e $995$}
\resolucao{(a)
\[
\mu_r=\frac{B}{\mu_0H}:\quad
\frac{1{,}2}{(4\pi\times10^{-7})(400)}=2390;
\qquad
\frac{1{,}5}{(4\pi\times10^{-7})(1200)}=995.
\]
(b) \textbf{A permeabilidade caiu para menos da metade} enquanto $H$
triplicou.\\
\textbf{Interpretação:} o material está entrando na região de \emph{joelho} da
curva de magnetização. Os domínios já estão em boa parte alinhados, e cada
acréscimo de $H$ produz cada vez menos acréscimo de $B$.\\
\textbf{Consequência prática:} triplicar a corrente aumentou o fluxo em apenas
$25\%$. Operar nessa região é ineficiente --- o cobre gasta muito para
render pouco.\\
\textbf{Lição de fundo:} $\mu_r$ de catálogo é sempre um valor \emph{para uma
condição específica}. Usá-lo fora dessa condição é um erro comum e silencioso.}
\end{questao}

\begin{questao}[2]{}
Explique a diferença entre a permeabilidade \textbf{normal}
($\mu=B/H$) e a permeabilidade \textbf{diferencial}
($\mu_d=\mathrm{d}B/\mathrm{d}H$), e indique quando cada uma se aplica.
\resposta{Normal para o ponto de operação; diferencial para pequenos sinais
sobrepostos}
\resolucao{\textbf{Permeabilidade normal:} razão entre $B$ e $H$ no ponto de
operação. Ela governa o \emph{fluxo total} e, portanto, a relutância e o
dimensionamento do núcleo.\\
\textbf{Permeabilidade diferencial:} inclinação da curva naquele ponto. Ela
governa a resposta a \emph{pequenas variações} sobrepostas --- e portanto a
indutância incremental.\\
\textbf{Onde diferem drasticamente:} perto da saturação, $B/H$ ainda é grande
(centenas), mas $\mathrm{d}B/\mathrm{d}H$ já caiu para quase $\mu_0$. Um
indutor com corrente contínua elevada pode ter indutância incremental dezenas
de vezes menor que a nominal.\\
\textbf{Consequência de projeto:} em filtros com componente contínua --- muito
comuns em fontes chaveadas --- é a permeabilidade \emph{diferencial} que
determina o desempenho, e é ela que precisa ser consultada no catálogo.}
\end{questao}

\begin{questao}[2]{}
Compare a temperatura de Curie do ferro ($\SI{770}{\celsius}$), do níquel
($\SI{358}{\celsius}$) e de uma ferrite típica
($\approx\SI{250}{\celsius}$), e explique o que ocorre acima dela.
\resposta{Acima de $T_C$ o material torna-se paramagnético e perde
praticamente toda a permeabilidade}
\resolucao{\textbf{O que ocorre acima de $T_C$.} A agitação térmica supera a
interação de troca que mantém os momentos atômicos alinhados dentro dos
domínios. A estrutura de domínios se desfaz, e o material passa de
\textbf{ferromagnético} a \textbf{paramagnético}:
\[
\mu_r:\ 10^{3}\text{--}10^{4}\ \longrightarrow\ \approx1.
\]
\textbf{A transição é abrupta} --- trata-se de uma transição de fase de segunda
ordem, e $\mu_r$ despenca em poucos graus em torno de $T_C$.\\
\textbf{Consequências práticas.}
\begin{itemize}
\item \textbf{Ferrites são as mais vulneráveis.} Com $T_C$ próxima de
$\SI{250}{\celsius}$, e temperatura de operação já em torno de
$\SI{100}{\celsius}$, a margem é pequena. Um transformador de fonte chaveada
que superaqueça perde indutância, a corrente dispara e ele se destrói --- fuga
térmica.
\item \textbf{Aplicação deliberada:} panelas de indução com fundo cuja $T_C$ é
próxima da temperatura desejada param de aquecer sozinhas ao atingi-la. É um
termostato sem sensor.
\item \textbf{Desmagnetização por aquecimento:} aquecer um ímã acima de
$T_C$ e resfriá-lo sem campo o deixa completamente desmagnetizado.
\end{itemize}}
\end{questao}

\begin{questao}[2]{}
Um ímã permanente de ferrite tem coeficiente de temperatura de
$\SI{-0.2}{\percent\per\celsius}$ para a remanência. Determine a variação de
$B_r$ entre $\SI{20}{\celsius}$ e $\SI{120}{\celsius}$.
\resposta{Redução de $20\%$}
\resolucao{
\[
\frac{\Delta B_r}{B_r}=(-0{,}002)(100)=-0{,}20=-20\%.
\]
\textbf{Reversível ou não?} Até certo limite, a perda é \textbf{reversível} --- o
ímã recupera a magnetização ao esfriar. Acima de uma temperatura crítica
(bem abaixo de $T_C$), parte da perda torna-se \textbf{irreversível}, porque
domínios se reorientam permanentemente.\\
\textbf{Comparação entre materiais:}
\begin{center}
\begin{tabular}{lcc}
\hline
Material & Coef. de $B_r$ & $T_{\max}$ típica\\
\hline
Ferrite & $\SI{-0.2}{\percent\per\celsius}$ & $\SI{250}{\celsius}$\\
Neodímio (NdFeB) & $\SI{-0.12}{\percent\per\celsius}$ & $\SI{80}{}$ a
 $\SI{200}{\celsius}$\\
Samário-cobalto & $\SI{-0.03}{\percent\per\celsius}$ & $\SI{300}{\celsius}$\\
\hline
\end{tabular}
\end{center}
\textbf{Compromisso revelador:} o neodímio tem a maior densidade de energia mas
péssima estabilidade térmica; o samário-cobalto é bem mais estável e mais caro.
Em motores que aquecem, a escolha raramente é o material mais forte.}
\end{questao}

\begin{questao}[2]{}
Um núcleo de aço-silício dissipa perdas por histerese proporcionais à área do
ciclo. Com área de $\SI{200}{\joule\per\meter\cubed}$ por ciclo e volume de
$\SI{1000}{\centi\meter\cubed}$, determine a potência a $\SI{60}{\hertz}$ e a
$\SI{400}{\hertz}$.
\resposta{$\SI{12}{\watt}$ e $\SI{80}{\watt}$}
\resolucao{
\[
P_h=f\times(\text{área})\times V:
\quad
60(200)(\pot{1}{-3})=\SI{12}{\watt};
\qquad
400(200)(\pot{1}{-3})=\SI{80}{\watt}.
\]
\textbf{Proporcionalidade linear com $f$} --- cada ciclo custa a mesma energia,
e mais ciclos por segundo significam mais potência.\\
\textbf{Contraste com as correntes parasitas}, cujas perdas vão com
$f^{2}$. Em baixa frequência a histerese domina; em alta, as parasitas.\\
\textbf{Consequência:} o mesmo núcleo que funciona bem em $\SI{60}{\hertz}$
pode ser inviável em $\SI{400}{\hertz}$ --- e é por isso que equipamentos
aeronáuticos, que operam nessa frequência, usam ligas especiais de lâminas
finíssimas.}
\end{questao}

\begin{questao}[2]{}
Um material magnético \textbf{mole} tem $H_c=\SI{50}{\ampere\per\meter}$ e um
material \textbf{duro} tem $H_c=\SI{500}{\kilo\ampere\per\meter}$. Compare-os
e indique aplicações.
\resposta{Razão de $10^{4}$ na coercividade; moles para núcleos, duros para ímãs
permanentes}
\resolucao{\textbf{Coercividade $H_c$} é o campo necessário para \emph{anular} a
magnetização remanente. Ela mede o quanto o material "resiste" a mudar de
estado.
\begin{center}
\begin{tabular}{lccl}
\hline
Tipo & $H_c$ & Área do ciclo & Aplicação\\
\hline
Mole & baixa & pequena & núcleos de transformador e motor\\
Duro & alta & grande & ímãs permanentes\\
\hline
\end{tabular}
\end{center}
\textbf{Materiais moles} (ferro-silício, permalloy, ferrites moles):
magnetizam e desmagnetizam facilmente, com pouca perda por ciclo. É o que se
quer em algo que inverte a magnetização $120$ vezes por segundo.\\
\textbf{Materiais duros} (alnico, ferrite dura, NdFeB): resistem à
desmagnetização, mantendo o campo sem corrente alguma.\\
\textbf{O critério é a área do ciclo:} pequena para quem vai percorrê-lo muitas
vezes, grande para quem precisa permanecer magnetizado.}
\end{questao}

\begin{questao}[2]{}
Explique por que a curva de magnetização \textbf{inicial} de um material
virgem difere do ciclo de histerese subsequente.
\resposta{A curva inicial parte da origem; o ciclo nunca mais passa por ela}
\resolucao{\textbf{Material virgem:} os domínios estão orientados
aleatoriamente, e a magnetização líquida é nula. Aplicando $H$ crescente, a
curva parte da \textbf{origem} e sobe até a saturação --- é a \emph{curva de
magnetização inicial}.\\
\textbf{Ao reduzir $H$ a zero}, a magnetização \emph{não} retorna a zero: resta
a \textbf{remanência} $B_r$. Os domínios que se reorientaram não voltam
espontaneamente.\\
\textbf{A partir daí}, ciclando $H$ entre $\pm H_{\max}$, o material percorre
um laço fechado que \textbf{nunca mais passa pela origem}.\\
\textbf{Como voltar ao estado virgem:} aplicar campo alternado de amplitude
decrescente até zero (desmagnetização por ciclos decrescentes), ou aquecer acima
de $T_C$ e resfriar sem campo.\\
\textbf{Importância prática:} ao levantar experimentalmente a curva
$B\times H$, é preciso desmagnetizar o material antes --- caso contrário, os
primeiros pontos refletem a história anterior, não a propriedade do material.}
\end{questao}

\begin{questao}[2]{}
Um transformador é desligado no pico da corrente de magnetização. Explique o
que ocorre na religação e por que a corrente de \emph{inrush} pode ser elevada.
\resposta{O fluxo remanente soma-se ao fluxo imposto, podendo levar o núcleo a
saturação profunda}
\resolucao{\textbf{Ao desligar}, o núcleo retém fluxo remanente
$\Phi_r$ --- consequência direta da histerese.\\
\textbf{Ao religar}, o fluxo evolui a partir desse valor. No pior caso, a tensão
é aplicada no instante em que exigiria fluxo de mesma polaridade, e o fluxo
total pode chegar a
\[
\Phi\approx\Phi_r+2\Phi_{\max},
\]
muito além da saturação.\\
\textbf{Consequência:} saturado, o núcleo perde permeabilidade, a indutância
desaba e a corrente é limitada apenas pela resistência do enrolamento. A
corrente de \emph{inrush} pode atingir dez a vinte vezes a nominal, por alguns
ciclos.\\
\textbf{Soluções empregadas:}
\begin{itemize}
\item \textbf{Chaveamento sincronizado:} energizar no instante correto da onda
de tensão.
\item \textbf{Resistor ou NTC de pré-carga}, curto-circuitado após alguns
ciclos.
\item \textbf{Proteção com curva temporizada}, que tolera o transitório sem
atuar.
\end{itemize}
\textbf{Note a conexão:} este fenômeno só existe \emph{porque} há histerese ---
um núcleo hipotético sem remanência não teria \emph{inrush} por essa causa.}
\end{questao}

\begin{questao}[2]{}
Um circuito magnético tem dois trechos em série: ferro
($\ell_1=\SI{30}{\centi\meter}$, $\mu_r=2000$) e ar
($\ell_2=\SI{2}{\milli\meter}$), ambos com seção
$\SI{5}{\centi\meter\squared}$. Determine a relutância total e a fração de cada
trecho.
\resposta{$\mathcal{R}_{fe}=\pot{2.39}{5}$, $\mathcal{R}_{ar}=\pot{3.18}{6}$;
o ar responde por $93\%$}
\resolucao{
\[
\mathcal{R}_{fe}=\frac{0{,}30}{(2000)\mu_0(\pot{5}{-4})}
=\pot{2.39}{5}\,\si{\per\henry};
\]
\[
\mathcal{R}_{ar}=\frac{\pot{2}{-3}}{\mu_0(\pot{5}{-4})}
=\pot{3.18}{6}\,\si{\per\henry}.
\]
\[
\mathcal{R}_{tot}=\pot{3.42}{6};
\qquad
\frac{\mathcal{R}_{ar}}{\mathcal{R}_{tot}}=93\%.
\]
\textbf{Dois milímetros de ar dominam trinta centímetros de ferro} --- porque a
razão de permeabilidades ($2000$) supera de longe a razão de comprimentos
($150$).\\
\textbf{As relutâncias somam-se em série}, exatamente como resistências. E o
fluxo é o mesmo nos dois trechos, o que garante $B$ igual --- já que a seção
também é a mesma.\\
\textbf{Consequência:} o comportamento do circuito é governado quase
inteiramente pelo entreferro, e a não linearidade do ferro fica mascarada. É
esse efeito que \emph{lineariza} indutores com entreferro.}
\end{questao}

\blocoaprofundamento

\begin{questao}[3]{}
Distinga rigorosamente $\vec H$, $\vec B$ e $\vec M$.
\begin{itensq}
\item Escreva a relação entre elas e as unidades.
\item Determine qual delas muda ao inserir um núcleo em um solenoide, mantida a
      corrente.
\item Explique qual é "mais fundamental" e por quê.
\end{itensq}
\resposta{$\vec B=\mu_0(\vec H+\vec M)$; ao inserir o núcleo, $H$ não muda e
$B$ aumenta}
\resolucao{(a)
\[
\vec B=\mu_0\left(\vec H+\vec M\right),
\]
com $[H]=[M]=\si{\ampere\per\meter}$ e $[B]=\si{\tesla}$.\\
\emph{($H$ é às vezes chamado "intensidade de campo magnético" e $B$
"densidade de fluxo" --- nomenclatura infeliz, porque sugere que $H$ é o campo
principal.)}\\
(b) \textbf{Ao inserir o núcleo com corrente constante:}
\begin{itemize}
\item $H=NI/\ell$ \textbf{não muda} --- ele depende só das correntes de
condução.
\item $M$ \textbf{passa de zero a $(\mu_r-1)H$} --- o material se magnetiza.
\item $B$ \textbf{aumenta por} $\mu_r$ --- consequência das duas anteriores.
\end{itemize}
(c) \textbf{Qual é mais fundamental: $\vec B$.} Argumentos:
\begin{itemize}
\item É $\vec B$ que aparece na força de Lorentz
($\vec F=q\vec v\times\vec B$) --- é ele que age sobre cargas.
\item É $\vec B$ que aparece na lei de Faraday --- é seu fluxo que induz f.e.m.
\item $\nabla\cdot\vec B=0$ é lei universal; $\nabla\cdot\vec H=-\nabla\cdot
\vec M$ não é nula em geral.
\end{itemize}
\textbf{Para que serve $\vec H$, então:} ele é conveniente porque sua circulação
depende \emph{apenas} das correntes de condução ($\oint\vec H\cdot
\mathrm{d}\vec\ell=NI$), sem precisar contabilizar as correntes de
magnetização do material. É uma grandeza \emph{auxiliar} de enorme utilidade
prática --- mas auxiliar.}
\end{questao}

\begin{questao}[3]{}
A curva $B\times H$ de um material apresenta três regiões distintas.
\begin{itensq}
\item Descreva cada uma em termos de domínios.
\item Identifique onde $\mu_r$ é máxima.
\item Indique a região de operação recomendada.
\end{itensq}
\resposta{Deslocamento reversível, deslocamento irreversível e rotação;
$\mu_r$ máxima no trecho intermediário}
\resolucao{(a) \textbf{Três regiões.}
\begin{enumerate}
\item \textbf{Campos baixos --- deslocamento reversível de paredes.} As paredes
entre domínios deslocam-se ligeiramente, favorecendo os domínios alinhados com
$H$. Removido o campo, elas retornam. A curva é pouco inclinada.
\item \textbf{Campos médios --- deslocamento irreversível (saltos de
Barkhausen).} As paredes superam obstáculos (impurezas, contornos de grão) em
saltos abruptos. É aqui que a curva é mais \textbf{íngreme} e onde nasce a
histerese, pois os saltos não se desfazem ao reduzir $H$.
\item \textbf{Campos altos --- rotação dos domínios.} Já não há paredes a
mover; os momentos giram progressivamente até se alinharem com $H$. A curva
achata-se: é o \emph{joelho} e depois a saturação.
\end{enumerate}
(b) \textbf{$\mu_r$ é máxima na região intermediária}, onde
$\mathrm{d}B/\mathrm{d}H$ é maior. Nem no início (paredes presas) nem no fim
(domínios já alinhados).\\
\textbf{Consequência contraintuitiva:} a permeabilidade em campos muito baixos é
\emph{menor} que a máxima --- daí catálogos especificarem $\mu_i$ (inicial) e
$\mu_{\max}$ separadamente.\\
(c) \textbf{Região recomendada:} abaixo do joelho, tipicamente com
$B_{\max}$ entre $60$ e $80\%$ de $B_{sat}$. Isso aproveita a alta
permeabilidade sem risco de saturação em transitórios.}
\end{questao}

\begin{questao}[3]{}
Um núcleo tem curva $B\times H$ não linear, e é preciso determinar o ponto de
operação com entreferro.
\begin{itensq}
\item Formule o problema.
\item Descreva o método gráfico da reta de entreferro.
\item Explique por que o entreferro estabiliza o ponto de operação.
\end{itensq}
\resposta{Interseção da curva $B\times H$ do material com a reta imposta pelo
entreferro}
\resolucao{(a) \textbf{Formulação.} A lei de Ampère no circuito magnético dá
\[
NI=H_{fe}\,\ell_{fe}+H_g\,g,
\]
e no entreferro $H_g=B/\mu_0$ (com $B$ igual ao do ferro, por conservação do
fluxo e seção constante). Então
\[
NI=H_{fe}\ell_{fe}+\frac{B}{\mu_0}g.
\]
Duas incógnitas ($B$ e $H_{fe}$) e uma equação --- a segunda relação é a
\emph{curva do material}, que não tem forma algébrica simples.\\
(b) \textbf{Método gráfico.} Reescrevendo,
\[
B=\frac{\mu_0\,NI}{g}-\frac{\mu_0\ell_{fe}}{g}\,H_{fe},
\]
que é uma \textbf{reta} no plano $B\times H$ --- a \emph{reta de entreferro} ou
\emph{reta de carga magnética}.\\
Traçando-a sobre a curva $B\times H$ do material, o \textbf{ponto de operação} é
a interseção.\\
\emph{(É exatamente o mesmo procedimento da reta de carga usada com diodos: uma
característica não linear do dispositivo e uma reta imposta pelo circuito.)}\\
(c) \textbf{Por que o entreferro estabiliza.} A inclinação da reta é
$-\mu_0\ell_{fe}/g$: quanto \emph{maior} o entreferro, mais \textbf{deitada} a
reta.\\
Uma reta deitada cruza a curva do material em uma região de inclinação
comparável --- e ali a interseção é \emph{pouco sensível} a variações de
$\mu_r$ do material. Trocar o lote de ferro, ou aquecê-lo, quase não desloca o
ponto.\\
\textbf{Sem entreferro}, a reta é quase vertical e a interseção depende
inteiramente da curva do material --- que varia com temperatura, lote e
história magnética. \textbf{O entreferro compra previsibilidade ao custo de mais
ampère-espiras.}}
\end{questao}

\begin{questao}[3]{}
Compare aço-silício e ferrite para aplicação em transformadores.
\begin{itensq}
\item Compare $B_{sat}$, resistividade e faixa de frequência.
\item Determine em qual faixa cada um é preferível.
\item Justifique fisicamente.
\end{itensq}
\resposta{Aço-silício: alto $B_{sat}$, baixa resistividade, até
$\si{\kilo\hertz}$; ferrite: baixo $B_{sat}$, altíssima resistividade,
até $\si{\mega\hertz}$}
\resolucao{(a)
\begin{center}
\begin{tabular}{lcc}
\hline
Propriedade & Aço-silício & Ferrite (MnZn)\\
\hline
$B_{sat}$ & $1{,}5$--$\SI{2}{\tesla}$ & $0{,}3$--$\SI{0.5}{\tesla}$\\
Resistividade & $\approx\pot{5}{-7}\,\si{\ohm\meter}$ &
 $10^{-1}$ a $\SI{10}{\ohm\meter}$\\
$\mu_r$ típica & $2000$--$8000$ & $1000$--$5000$\\
Faixa útil & até $\approx\SI{1}{\kilo\hertz}$ & $\SI{20}{\kilo\hertz}$ a
 $\si{\mega\hertz}$\\
\hline
\end{tabular}
\end{center}
\textbf{Note a resistividade:} a ferrite é um \emph{cerâmico}, com
resistividade até $10^{7}$ vezes maior que a do aço.\\
(b) \textbf{Baixa frequência} ($\SI{50}{}$--$\SI{400}{\hertz}$): aço-silício.
Seu $B_{sat}$ elevado permite núcleos menores para a mesma potência, e as perdas
parasitas são controladas por laminação.\\
\textbf{Alta frequência} ($>\SI{20}{\kilo\hertz}$): ferrite. Nenhuma laminação
seria fina o bastante para conter as correntes parasitas no aço.\\
(c) \textbf{Justificativa física.} As perdas por correntes parasitas vão com
\[
P_F\propto\frac{B^{2}f^{2}t^{2}}{\rho}.
\]
Em alta frequência, o termo $f^{2}$ explode. Há duas formas de compensar:
reduzir $t$ (laminação) ou aumentar $\rho$.\\
A ferrite ataca $\rho$ --- e ganha por muitas ordens de grandeza, o que nenhuma
laminação alcança.\\
\textbf{O preço:} $B_{sat}$ baixo obriga a usar seção maior ou aceitar menos
fluxo. \textbf{Mas em alta frequência isso não importa}, porque
$V=4{,}44Nf\Phi_{\max}$ --- com $f$ mil vezes maior, o fluxo necessário é mil
vezes menor. \textbf{É a alta frequência que torna a ferrite viável, e a ferrite
que torna a alta frequência viável.}}
\end{questao}

\begin{questao}[3]{}
Um ímã permanente é retirado de seu circuito magnético e deixado com os polos
livres ao ar.
\begin{itensq}
\item Explique o que ocorre com seu ponto de operação.
\item Determine o efeito de aproximar uma peça de ferro.
\item Justifique a prática de armazenar ímãs com "quebra-fluxo".
\end{itensq}
\resposta{Ao ar, o ponto de operação desce na curva de desmagnetização; o ferro
o eleva}
\resolucao{(a) \textbf{Ao ar livre.} O ímã precisa "empurrar" o fluxo por um
caminho de alta relutância --- o próprio ar em torno dele. Isso equivale a um
entreferro enorme, e a reta de carga fica muito deitada.\\
O ponto de operação desce na \emph{curva de desmagnetização} (segundo quadrante
do laço), afastando-se de $B_r$. O ímã fica submetido a campo desmagnetizante
produzido por ele mesmo.\\
(b) \textbf{Aproximando ferro.} O ferro oferece caminho de baixa relutância, e a
reta de carga se levanta. O ponto de operação sobe, aproximando-se de
$B_r$ --- o ímã fica "mais forte".\\
É por isso que um ímã em ferradura fechado por uma armadura de ferro sustenta
muito mais peso que aberto.\\
(c) \textbf{Quebra-fluxo (\emph{keeper}).} É uma peça de ferro doce colocada
entre os polos durante o armazenamento. Ela:
\begin{itemize}
\item mantém o ponto de operação alto, evitando desmagnetização progressiva;
\item protege contra campos externos adversos;
\item reduz a atração de limalha e detritos.
\end{itemize}
\textbf{Ressalva histórica:} a prática era essencial com alnico, cuja curva de
desmagnetização é frágil. Com neodímio moderno --- cuja curva é praticamente
reta no segundo quadrante --- o efeito é bem menor, embora o quebra-fluxo
continue útil por outras razões.}
\end{questao}

\begin{questao}[3]{}
Um material tem $B_r=\SI{1.2}{\tesla}$ e $H_c=\SI{900}{\kilo\ampere\per\meter}$.
\begin{itensq}
\item Estime o produto de energia máximo.
\item Compare com ferrite ($B_r=\SI{0.4}{\tesla}$,
      $H_c=\SI{250}{\kilo\ampere\per\meter}$).
\item Explique o significado do produto de energia.
\end{itensq}
\resposta{$(BH)_{\max}\approx\SI{270}{\kilo\joule\per\meter\cubed}$ contra
$\SI{25}{\kilo\joule\per\meter\cubed}$}
\resolucao{(a) Para curva de desmagnetização aproximadamente \emph{reta} (caso
do neodímio), o produto máximo ocorre no ponto médio:
\[
(BH)_{\max}\approx\frac{B_rH_c}{4}
=\frac{(1{,}2)(\pot{9}{5})}{4}
=\SI{270}{\kilo\joule\per\meter\cubed}.
\]
(b) Para a ferrite:
\[
(BH)_{\max}\approx\frac{(0{,}4)(\pot{2.5}{5})}{4}
=\SI{25}{\kilo\joule\per\meter\cubed},
\]
cerca de \textbf{onze vezes menor}.\\
(c) \textbf{Significado.} O produto de energia mede a energia magnética
disponível \emph{por unidade de volume} do ímã. Ele é a figura de mérito que
determina quão pequeno o ímã pode ser para produzir dado campo em dado
entreferro.\\
\textbf{Consequência prática:} para o mesmo desempenho, um ímã de neodímio ocupa
cerca de um décimo do volume de um de ferrite. Isso é o que viabilizou
servomotores compactos, fones de ouvido pequenos e discos rígidos.\\
\textbf{Contrapartida:} o neodímio é caro, oxida (precisa de revestimento),
desmagnetiza acima de $\SI{80}{}$--$\SI{200}{\celsius}$ e depende de
terras-raras. A ferrite é barata, quimicamente estável e térmicamente mais
tolerante --- e por isso continua dominando em volume de produção.}
\end{questao}

\begin{questao}[3]{}
Explique a diferença entre desmagnetização por \textbf{aquecimento} e por
\textbf{campo alternado decrescente}.
\begin{itensq}
\item Descreva cada processo.
\item Compare eficácia e aplicabilidade.
\end{itensq}
\resposta{Aquecimento acima de $T_C$ desordena os domínios; campo alternado os
distribui simetricamente}
\resolucao{(a) \textbf{Aquecimento acima de $T_C$.} A agitação térmica destrói o
alinhamento dentro dos domínios. Resfriando \emph{sem campo aplicado}, os
domínios se reformam em orientações aleatórias, e a magnetização líquida é
nula.\\
\textbf{Campo alternado decrescente.} Aplica-se campo alternado de amplitude
inicialmente alta (o bastante para saturar) e reduz-se lentamente a zero. O
material percorre laços de histerese cada vez menores, espiralando até a origem.
Ao final, tantos domínios apontam em um sentido quanto no outro.\\
(b) \textbf{Comparação.}
\begin{center}
\begin{tabular}{lcc}
\hline
& Aquecimento & Campo alternado\\
\hline
Eficácia & completa & muito boa\\
Danifica a peça? & possivelmente & não\\
Aplicável a peça montada? & raramente & sim\\
Rapidez & lento & segundos\\
\hline
\end{tabular}
\end{center}
\textbf{Na prática industrial, usa-se campo alternado} --- em desmagnetizadores
de ferramentas, de peças usinadas e de fitas magnéticas. O aquecimento é
inviável para peças que não toleram $\SI{770}{\celsius}$.\\
\textbf{Cuidado essencial:} a redução do campo deve ser \emph{lenta} em relação
ao período. Desligar bruscamente o desmagnetizador pode deixar a peça
magnetizada no valor correspondente à última meia-onda --- pior que antes.\\
\textbf{Alternativa equivalente:} afastar lentamente a peça de um campo
alternado fixo, o que produz o mesmo decaimento gradual.}
\end{questao}

\begin{questao}[3]{}
Um núcleo opera com $B_{\max}=\SI{1.2}{\tesla}$ a $\SI{60}{\hertz}$. Analise o
efeito de elevar $B_{\max}$ para $\SI{1.5}{\tesla}$.
\begin{itensq}
\item Determine o efeito sobre o número de espiras.
\item Determine o efeito sobre as perdas.
\item Avalie o compromisso.
\end{itensq}
\resposta{$N$ cai $20\%$; perdas por histerese sobem cerca de $60\%$ e
parasitas $56\%$}
\resolucao{(a) De $V=4{,}44Nf\Phi_{\max}=4{,}44NfB_{\max}A$, com $V$, $f$ e
$A$ fixos:
\[
N\propto\frac{1}{B_{\max}}
\;\Longrightarrow\;
\frac{N'}{N}=\frac{1{,}2}{1{,}5}=0{,}80.
\]
\textbf{Vinte por cento menos espiras} --- menos cobre, menos resistência, menos
perdas no enrolamento.\\
(b) \textbf{Perdas no núcleo aumentam.}
\begin{itemize}
\item \textbf{Parasitas:} $P_F\propto B_{\max}^{2}$, logo
$(1{,}5/1{,}2)^{2}=1{,}56$ --- aumento de $56\%$.
\item \textbf{Histerese:} $P_h\propto B_{\max}^{n}$ com $n\approx1{,}6$ a
$2$ (expoente de Steinmetz), dando aumento de $60$ a $56\%$.
\end{itemize}
(c) \textbf{O compromisso.}
\begin{center}
\begin{tabular}{lc}
\hline
Elevar $B_{\max}$ & Efeito\\
\hline
Espiras & $-20\%$\\
Perdas no cobre & diminuem\\
Perdas no núcleo & $+\approx58\%$\\
Volume e custo & diminuem\\
Margem para saturação & \textbf{diminui}\\
\hline
\end{tabular}
\end{center}
\textbf{Existe um ótimo econômico}, em que a soma das perdas é mínima --- e ele
depende do custo relativo do cobre e do aço e do regime de carga.\\
\textbf{Mas há um limite duro:} $\SI{1.5}{\tesla}$ já está próximo do joelho. A
margem para sobretensões e para o transitório de energização encolhe, e o risco
de \emph{inrush} severo cresce. \textbf{Projetar perto da saturação economiza
material e compra problemas operacionais.}}
\end{questao}

\begin{questao}[3]{}
Explique por que a permeabilidade de um material ferromagnético não é uma
constante, e quais parâmetros a afetam.
\resposta{$\mu$ depende de $H$, da temperatura, da frequência e da história
magnética}
\resolucao{\textbf{Quatro dependências, todas relevantes.}
\begin{enumerate}
\item \textbf{Da intensidade $H$.} É a mais óbvia: $\mu_r$ parte de um valor
inicial, cresce até um máximo e despenca na saturação. Variação de uma ordem de
grandeza ou mais.
\item \textbf{Da temperatura.} $\mu_r$ tende a \emph{crescer} ao se aproximar de
$T_C$ (efeito Hopkinson) e depois colapsa abruptamente. Isso torna o
comportamento térmico não monotônico e difícil de prever.
\item \textbf{Da frequência.} Em alta frequência, as paredes de domínio não
conseguem acompanhar o campo --- há \emph{relaxação}. A permeabilidade efetiva
cai, e aparece uma componente de perda. Catálogos de ferrite trazem
$\mu$ complexo, com parte real e imaginária.
\item \textbf{Da história magnética.} Por causa da histerese, o mesmo $H$ pode
corresponder a diferentes $B$, conforme o material esteja subindo ou descendo o
laço.
\end{enumerate}
\textbf{Consequência para o projetista.} Um valor único de $\mu_r$ em catálogo é
sempre uma \emph{simplificação} referida a condições específicas --- geralmente
campo baixo, temperatura ambiente e baixa frequência.\\
\textbf{Prática correta:} usar $\mu_r$ para estimativas iniciais, e recorrer às
\emph{curvas} do fabricante (de $B\times H$, de perdas por frequência, de
$\mu$ por temperatura) para o dimensionamento final. Projetos de fontes
chaveadas que ignoram isso costumam funcionar na bancada e falhar em campo.}
\end{questao}

\begin{questao}[3]{}
Um sensor de proximidade indutivo detecta metais pela variação de indutância.
\begin{itensq}
\item Explique o princípio para metais ferromagnéticos.
\item Explique para metais não ferromagnéticos (alumínio, cobre).
\item Determine por que o alcance difere entre eles.
\end{itensq}
\resposta{Ferromagnéticos aumentam $L$ por permeabilidade; não ferromagnéticos
reduzem $L$ e o fator de qualidade por correntes parasitas}
\resolucao{(a) \textbf{Ferromagnéticos.} A aproximação do material oferece
caminho de menor relutância, aumentando a indutância da bobina sensora. A
variação de $L$ é detectada pelo circuito.\\
Simultaneamente, correntes parasitas no aço amortecem o oscilador.\\
(b) \textbf{Não ferromagnéticos.} Alumínio e cobre têm $\mu_r\approx1$ --- não
alteram a relutância. Mas são \emph{bons condutores}, e o campo alternado induz
neles correntes parasitas intensas.\\
Pela lei de Lenz, essas correntes criam campo oposto, o que \textbf{reduz} a
indutância efetiva e, sobretudo, \textbf{drena energia} do oscilador, reduzindo
seu fator de qualidade.\\
(c) \textbf{Por que o alcance difere.} Sensores indutivos comuns são
especificados para aço, e apresentam \emph{fatores de correção}:
\begin{center}
\begin{tabular}{lc}
\hline
Material & Fator de alcance\\
\hline
Aço & $1{,}00$\\
Aço inoxidável & $\approx0{,}70$\\
Latão & $\approx0{,}50$\\
Alumínio & $\approx0{,}40$\\
Cobre & $\approx0{,}30$\\
\hline
\end{tabular}
\end{center}
\textbf{Razão:} no aço, os dois efeitos (permeabilidade e parasitas) somam-se.
Nos não ferromagnéticos, resta apenas o segundo --- e ele é tanto mais fraco
quanto \emph{melhor} o condutor, porque as correntes se concentram numa
película superficial cada vez mais fina.\\
\textbf{Consequência de aplicação:} um sensor ajustado para detectar peças de
aço a $\SI{10}{\milli\meter}$ só detectará cobre a $\SI{3}{\milli\meter}$ ---
erro de projeto frequente em linhas de montagem com peças mistas.}
\end{questao}

\blocodesafio

\begin{questao}[4]{}
\textbf{(Energia do ciclo de histerese.)} Demonstre que a área do laço
$B\times H$ representa energia dissipada por unidade de volume e por ciclo.
\begin{itensq}
\item Deduza a expressão.
\item Determine a potência dissipada em função da frequência.
\item Explique a origem microscópica da perda.
\end{itensq}
\resposta{$w=\oint H\,\mathrm{d}B$ por ciclo; $P=fVw$}
\resolucao{(a) \textbf{Dedução.} A potência instantânea fornecida ao núcleo pela
fonte é $p=\varepsilon i$. Com $\varepsilon=NA\,\mathrm{d}B/\mathrm{d}t$ e
$i=H\ell/N$:
\[
p=NA\frac{\mathrm{d}B}{\mathrm{d}t}\cdot\frac{H\ell}{N}
=(A\ell)\,H\,\frac{\mathrm{d}B}{\mathrm{d}t}
=V\,H\,\frac{\mathrm{d}B}{\mathrm{d}t}.
\]
A energia por ciclo é
\[
W=\int p\,\mathrm{d}t=V\oint H\,\mathrm{d}B,
\]
e $\oint H\,\mathrm{d}B$ é precisamente a \textbf{área do laço} no plano
$B\times H$, com unidade $\si{\joule\per\meter\cubed}$.\\
\textbf{Se o material fosse sem histerese}, o laço colapsaria em uma curva
única, a integral de ida cancelaria a de volta e a energia por ciclo seria
\textbf{zero} --- toda a energia magnética seria devolvida.\\
(b)
\[
P_h=f\,V\oint H\,\mathrm{d}B,
\]
\textbf{linear em $f$}. (As perdas por correntes parasitas, distintas destas,
vão com $f^{2}$.)\\
(c) \textbf{Origem microscópica.} A perda vem dos \textbf{saltos de Barkhausen}:
as paredes de domínio ficam presas em impurezas, defeitos e contornos de grão, e
se libertam abruptamente quando o campo cresce o suficiente.\\
Cada salto é um processo \textbf{irreversível} --- a energia elástica acumulada
converte-se em vibração da rede, ou seja, em \textbf{calor}.\\
\textbf{Consequência de engenharia:} materiais mais puros e com grãos maiores e
melhor orientados têm menos obstáculos, laços mais estreitos e menores perdas.
É exatamente isso que o processo de \emph{grão orientado} do aço-silício
persegue.}
\end{questao}

\begin{questao}[4]{}
\textbf{(Circuito magnético com dois materiais.)} Um circuito tem trecho de
ferro ($\ell_1$, $\mu_1$) em série com trecho de ferrite ($\ell_2$,
$\mu_2$), mesma seção.
\begin{itensq}
\item Formule as equações.
\item Determine $B$ e $H$ em cada trecho.
\item Explique qual grandeza é contínua na interface e por quê.
\end{itensq}
\resposta{$B$ é contínuo; $H$ é descontínuo, na razão inversa das
permeabilidades}
\resolucao{(a) \textbf{Duas equações.}\\
\emph{Conservação do fluxo} (equivalente à LKC):
\[
\Phi_1=\Phi_2
\;\Longrightarrow\;
B_1A=B_2A
\;\Longrightarrow\;
B_1=B_2=B.
\]
\emph{Lei de Ampère} (equivalente à LKT):
\[
NI=H_1\ell_1+H_2\ell_2.
\]
(b) Com $H_i=B/\mu_i$:
\[
NI=B\left(\frac{\ell_1}{\mu_1}+\frac{\ell_2}{\mu_2}\right)
\;\Longrightarrow\;
B=\frac{NI}{\dfrac{\ell_1}{\mu_1}+\dfrac{\ell_2}{\mu_2}}.
\]
\[
H_1=\frac{B}{\mu_1};
\qquad
H_2=\frac{B}{\mu_2};
\qquad
\frac{H_1}{H_2}=\frac{\mu_2}{\mu_1}.
\]
(c) \textbf{Qual grandeza é contínua --- e por quê.}
\begin{itemize}
\item \textbf{$B$ é contínuo} através da interface (na componente normal),
porque $\nabla\cdot\vec B=0$: o fluxo que chega precisa continuar. Não há
"acúmulo de fluxo" possível.
\item \textbf{$H$ é descontínuo}, na razão inversa das permeabilidades. O
material de menor $\mu$ exige \emph{muito mais} $H$ para sustentar o mesmo
$B$.
\end{itemize}
\textbf{Exemplo numérico revelador.} Com $\mu_{r1}=2000$ (ferro) e
$\mu_{r2}=2$ (um material quase não magnético), $H_2$ seria \emph{mil vezes}
$H_1$ --- e praticamente toda a força magnetomotriz seria consumida naquele
trecho, ainda que ele fosse curtíssimo.\\
\textbf{É exatamente o que acontece no entreferro} ($\mu_r=1$), e por isso ele
domina o circuito. Um entreferro não é um caso especial: é apenas o caso extremo
deste problema de dois materiais.}
\end{questao}

\begin{questao}[4]{}
\textbf{(Ensaio do ciclo de histerese.)} Projete um experimento para levantar o
laço $B\times H$ de um material em osciloscópio.
\begin{itensq}
\item Descreva o circuito e o que cada canal mede.
\item Determine as constantes de conversão.
\item Aponte os cuidados experimentais.
\end{itensq}
\resposta{Canal horizontal $\propto H$ (corrente por shunt); vertical
$\propto B$ (f.e.m. integrada)}
\resolucao{(a) \textbf{Circuito.} Um núcleo toroidal com dois enrolamentos:
primário $N_1$ excitado por fonte alternada e secundário $N_2$ para medição.
\begin{itemize}
\item \textbf{Canal horizontal ($X$):} tensão sobre um resistor \emph{shunt}
$R_s$ em série com o primário. Como $H=N_1i_1/\ell$, essa tensão é proporcional
a $H$.
\item \textbf{Canal vertical ($Y$):} a f.e.m. do secundário passa por um
\textbf{integrador} $RC$. Como
$\varepsilon_2=N_2A\,\mathrm{d}B/\mathrm{d}t$, sua integral é proporcional a
$B$.
\end{itemize}
Em modo $XY$, o osciloscópio traça diretamente o laço.\\
(b) \textbf{Constantes de conversão.}
\[
H=\frac{N_1}{\ell R_s}\,v_X;
\qquad
B=\frac{R_iC_i}{N_2A}\,v_Y,
\]
com $R_iC_i$ os componentes do integrador.\\
\textbf{Condição do integrador:} $R_iC_i\gg1/f$, para que ele integre de fato em
vez de apenas atenuar. Com $f=\SI{60}{\hertz}$, adote
$R_iC_i\ge\SI{0.1}{\second}$.\\
(c) \textbf{Cuidados experimentais.}
\begin{itemize}
\item \textbf{Desmagnetizar antes} --- caso contrário os primeiros ciclos
refletem a história anterior. Comece com amplitude alta e reduza.
\item \textbf{$R_s$ pequeno} (fração de ohm), para não distorcer a corrente nem
dissipar demais. Use quatro fios se possível.
\item \textbf{Núcleo toroidal}, que garante $H$ uniforme e caminho médio bem
definido --- em núcleos em E o caminho é ambíguo.
\item \textbf{Deriva do integrador:} qualquer \emph{offset} faz o laço "subir"
lentamente na tela. Use acoplamento CA ou integrador com realimentação CC.
\item \textbf{Varrer a amplitude} para obter a família de laços e localizar a
curva de magnetização normal, que é o lugar geométrico dos vértices.
\end{itemize}
\textbf{O que se extrai:} $B_r$ (interseção com o eixo vertical), $H_c$
(interseção com o horizontal), $B_{sat}$ e a área --- e daí as perdas por
ciclo.}
\end{questao}

\begin{questao}[4]{}
\textbf{(Domínios magnéticos.)} Explique a estrutura de domínios e sua relação
com as propriedades macroscópicas.
\begin{itensq}
\item Explique por que os domínios existem.
\item Relacione a estrutura com $B_r$, $H_c$ e $\mu_r$.
\item Explique a diferença entre materiais moles e duros nesse nível.
\end{itensq}
\resposta{Domínios minimizam a energia magnetostática; a mobilidade das paredes
determina $\mu_r$ e $H_c$}
\resolucao{(a) \textbf{Por que existem.} Um cristal ferromagnético
uniformemente magnetizado criaria campo externo intenso, com alta energia
magnetostática armazenada no espaço.\\
Dividir-se em \textbf{domínios} de orientações opostas faz esse campo externo
quase desaparecer --- as linhas se fecham internamente. O custo é a energia das
\emph{paredes de domínio} (paredes de Bloch), onde os momentos giram
gradualmente.\\
\textbf{O equilíbrio} entre esses dois termos fixa o tamanho típico dos
domínios, tipicamente de micrômetros a dezenas de micrômetros.\\
(b) \textbf{Relação com as propriedades macroscópicas.}
\begin{itemize}
\item \textbf{$\mu_r$ alta} exige paredes \emph{móveis} --- material puro, sem
tensões internas, com grãos grandes e bem orientados.
\item \textbf{$H_c$ baixa} tem a mesma origem: paredes que deslizam facilmente
não resistem à inversão.
\item \textbf{$B_r$ alta} exige que os domínios permaneçam alinhados após a
retirada do campo --- o que requer \emph{anisotropia} cristalina forte,
"ancorando" a magnetização em direções preferenciais.
\end{itemize}
(c) \textbf{Moles contra duros, no nível microscópico.}
\begin{center}
\begin{tabular}{lll}
\hline
& Moles & Duros\\
\hline
Paredes & muito móveis & fixadas ou inexistentes\\
Impurezas & minimizadas & introduzidas de propósito\\
Anisotropia & baixa & alta\\
Grãos & grandes, orientados & finos, monodomínio\\
\hline
\end{tabular}
\end{center}
\textbf{O contraste é elegante:} ímãs modernos de neodímio são feitos de
partículas tão pequenas que cada uma é um \textbf{monodomínio} --- não há
paredes para mover, e inverter a magnetização exige girar todos os momentos
contra a anisotropia. Daí a coercividade enorme.\\
\textbf{Já o aço-silício de grão orientado} persegue o oposto: grãos grandes,
alinhados por laminação a frio, com paredes livres para deslizar.\\
\textbf{Mesma física, objetivos opostos --- e processos de fabricação
inteiramente diferentes.}}
\end{questao}

\begin{questao}[4]{}
\textbf{(Projeto sob restrição de perdas.)} Um transformador de
$\SI{60}{\hertz}$ deve ter perdas no núcleo inferiores a $\SI{2}{\watt}$, com
material cuja perda específica é
$\SI{1.5}{\watt\per\kilogram}$ a $\SI{1.5}{\tesla}$ e
$\SI{0.9}{\watt\per\kilogram}$ a $\SI{1.2}{\tesla}$.
\begin{itensq}
\item Determine a massa admissível em cada condição.
\item Determine o efeito sobre o número de espiras.
\item Estabeleça o critério de escolha.
\end{itensq}
\resposta{$\SI{1.33}{\kilogram}$ a $\SI{1.5}{\tesla}$;
$\SI{2.22}{\kilogram}$ a $\SI{1.2}{\tesla}$}
\resolucao{(a)
\[
m\le\frac{2}{1{,}5}=\SI{1.33}{\kilogram}\ \text{a } \SI{1.5}{\tesla};
\qquad
m\le\frac{2}{0{,}9}=\SI{2.22}{\kilogram}\ \text{a } \SI{1.2}{\tesla}.
\]
(b) De $V=4{,}44NfB_{\max}A$, com tensão e frequência fixas,
\[
N\,B_{\max}A=\text{const}.
\]
Reduzir $B_{\max}$ de $1{,}5$ para $\SI{1.2}{\tesla}$ exige aumentar
$NA$ em $25\%$ --- mais espiras, ou seção maior, ou ambos.\\
(c) \textbf{Critério de escolha.} As duas opções levam a projetos diferentes:
\begin{center}
\begin{tabular}{lcc}
\hline
& $B=\SI{1.5}{\tesla}$ & $B=\SI{1.2}{\tesla}$\\
\hline
Massa de núcleo & menor & maior ($+67\%$)\\
Espiras & menos & mais ($+25\%$)\\
Perdas no cobre & menores & maiores\\
Margem para saturação & pequena & confortável\\
Corrente de \emph{inrush} & severa & moderada\\
\hline
\end{tabular}
\end{center}
\textbf{Note o conflito:} operar em $B$ alto economiza \emph{ferro} mas gasta
mais \emph{cobre} por unidade de perda; operar em $B$ baixo faz o inverso.\\
\textbf{Decisão:}
\begin{itemize}
\item \textbf{Se o custo de material dominar} (transformadores de baixo custo,
uso intermitente), escolha $B$ alto.
\item \textbf{Se o custo de energia dominar} (operação contínua por décadas),
escolha $B$ baixo --- as perdas em vazio existem $\SI{8760}{\hour}$ por ano.
\item \textbf{Se houver risco de sobretensão}, escolha $B$ baixo por segurança:
um transformador operando a $\SI{1.5}{\tesla}$ satura com apenas $10\%$ de
sobretensão.
\end{itemize}}
\end{questao}

\begin{questao}[4]{}
\textbf{(Anisotropia e grão orientado.)} Explique o que é o aço-silício de grão
orientado e por que ele é usado em transformadores mas não em motores.
\begin{itensq}
\item Explique o processo e o efeito.
\item Compare as perdas nas duas direções.
\item Justifique a diferença de aplicação.
\end{itensq}
\resposta{Grãos alinhados numa direção preferencial; excelente ao longo dela,
ruim transversalmente}
\resolucao{(a) \textbf{Processo e efeito.} O aço-silício é laminado a frio e
recozido de modo que os grãos cristalinos se alinhem com o eixo de fácil
magnetização do ferro (a direção $\langle100\rangle$ da célula cúbica) paralelo
à direção de laminação.\\
\textbf{Resultado:} nessa direção, magnetizar é muito mais fácil --- as paredes
de domínio se movem livremente, $\mu_r$ é altíssima e as perdas caem
drasticamente.\\
(b) \textbf{Comparação por direção.}
\begin{center}
\begin{tabular}{lcc}
\hline
& Direção de laminação & Transversal\\
\hline
Perdas específicas & $\approx\SI{1}{\watt\per\kilogram}$ &
 $2$ a $3\times$ maiores\\
$\mu_r$ & muito alta & bem menor\\
$B$ para dado $H$ & alto & reduzido\\
\hline
\end{tabular}
\end{center}
\textbf{A anisotropia é grande} --- o material é excelente em uma direção e
medíocre nas outras.\\
(c) \textbf{Por que serve a transformadores e não a motores.}
\begin{itemize}
\item \textbf{Transformador:} o fluxo percorre um caminho \emph{fixo e
conhecido} --- as colunas e as culatras do núcleo. É possível cortar as lâminas
de modo que o fluxo siga sempre a direção favorável. \emph{(Nos cantos, onde o
fluxo precisa dobrar, usam-se juntas em $\ang{45}$ para minimizar a
transição.)}
\item \textbf{Motor:} o fluxo no estator \textbf{gira}, percorrendo todas as
direções ao longo de cada revolução. Não existe direção preferencial a
privilegiar --- metade do tempo o fluxo estaria na direção ruim.
\end{itemize}
\textbf{Por isso motores usam aço de grão \emph{não} orientado}, deliberadamente
\textbf{isotrópico}: pior que o grão orientado em sua melhor direção, mas
uniforme em todas --- e portanto melhor \emph{na média} sobre uma rotação
completa.\\
\textbf{Princípio geral:} otimizar para o caso médio quando a solicitação é
variável, e para o caso específico quando ela é fixa.}
\end{questao}

\begin{questao}[4]{}
\textbf{(Blindagem magnética com material de alta permeabilidade.)} Analise
quantitativamente a blindagem de campo estático por invólucro de mu-metal.
\begin{itensq}
\item Explique o mecanismo por relutância.
\item Estime o fator de blindagem para invólucro cilíndrico.
\item Aponte as limitações práticas.
\end{itensq}
\resposta{Fator $\approx\mu_r t/D$; o mecanismo é desvio de fluxo, não
cancelamento}
\resolucao{(a) \textbf{Mecanismo.} Ao contrário da blindagem elétrica --- que
\emph{cancela} o campo por redistribuição de cargas --- a blindagem magnética
\textbf{desvia} o fluxo.\\
O invólucro de alta $\mu_r$ oferece caminho de relutância muito menor que o ar
interno. O fluxo "prefere" circular pela parede e contorna a região protegida.\\
\textbf{Analogia:} é um divisor de fluxo, em que o material é o ramo de baixa
relutância e a cavidade, o de alta.\\
(b) \textbf{Fator de blindagem.} Para um cilindro de diâmetro $D$ e parede
$t$, com $t\ll D$, a estimativa clássica é
\[
S\approx\frac{\mu_r\,t}{D}.
\]
Com $\mu_r=20000$ (mu-metal recozido), $t=\SI{1}{\milli\meter}$ e
$D=\SI{100}{\milli\meter}$:
\[
S\approx\frac{20000(1)}{100}=200,
\]
ou seja, atenuação de cerca de $\SI{46}{\decibel}$.\\
\textbf{Compare com a blindagem elétrica}, que atinge muitas ordens de grandeza
com uma folha fina de qualquer metal. A magnética é \emph{muito} mais modesta.\\
(c) \textbf{Limitações práticas.}
\begin{itemize}
\item \textbf{Saturação.} Se o campo externo for intenso, a parede satura,
$\mu_r$ desaba e a blindagem falha justamente quando mais se precisa dela.
\emph{Solução:} camadas múltiplas, com a externa de material de alto
$B_{sat}$ (aço comum) e a interna de mu-metal.
\item \textbf{Sensibilidade mecânica.} O mu-metal perde permeabilidade se for
dobrado, cortado ou submetido a choque. Exige \textbf{recozimento após a
conformação}, em atmosfera de hidrogênio --- processo caro.
\item \textbf{Aberturas.} Furos e fendas quebram o caminho de baixa relutância e
degradam a blindagem muito mais que no caso elétrico.
\item \textbf{Custo e peso}, ambos elevados.
\end{itemize}
\textbf{Conclusão:} blindar campo magnético estático é caro, frágil e de eficácia
limitada. Sempre que possível, é preferível \emph{afastar} a fonte, reduzir a
área das espiras ou compensar ativamente o campo com bobinas.}
\end{questao}

\begin{questao}[4]{}
\textbf{(Ponto de operação de ímã permanente.)} Determine o ponto de operação
de um ímã em um circuito magnético com entreferro.
\begin{itensq}
\item Formule a equação da reta de carga.
\item Determine a condição de máximo aproveitamento.
\item Explique o efeito de aumentar o entreferro.
\end{itensq}
\resposta{A reta de carga tem inclinação
$-\dfrac{\ell_m A_g}{g A_m}$; o aproveitamento é máximo em
$(BH)_{\max}$}
\resolucao{(a) \textbf{Reta de carga.} Duas condições:
\emph{conservação do fluxo} ($B_mA_m=B_gA_g$) e \emph{lei de Ampère}
($H_m\ell_m+H_gg=0$, pois não há corrente).\\
Da segunda, $H_m\ell_m=-H_gg=-\dfrac{B_g}{\mu_0}g$. Combinando com a primeira:
\[
B_m=-\mu_0\,\frac{\ell_m A_g}{g\,A_m}\,H_m.
\]
É uma reta pela \textbf{origem} no segundo quadrante ($H_m<0$, $B_m>0$) ---
a chamada \emph{linha de carga} do ímã.\\
\textbf{Note:} o ímã opera com $H$ \emph{negativo} --- ele está sujeito ao
próprio campo desmagnetizante. Sem entreferro ($g\to0$), a reta seria vertical e
o ponto de operação estaria em $B_r$.\\
(b) \textbf{Máximo aproveitamento.} A energia disponível no entreferro é
proporcional ao produto $B_mH_m$. O ponto ótimo é aquele em que
$|B_mH_m|$ é \textbf{máximo} --- o já mencionado $(BH)_{\max}$.\\
Dimensiona-se o circuito \emph{para que a reta de carga passe por esse ponto},
escolhendo a razão $\dfrac{\ell_m A_g}{g A_m}$ adequada. Isso minimiza o volume
de ímã necessário.\\
(c) \textbf{Efeito de aumentar o entreferro.} A inclinação
$-\mu_0\ell_mA_g/(gA_m)$ diminui em módulo --- a reta \textbf{deita}. O ponto de
operação desce na curva de desmagnetização:
\begin{itemize}
\item $B_m$ diminui (menos fluxo);
\item $|H_m|$ aumenta (mais campo desmagnetizante sobre o próprio ímã);
\item se o entreferro for grande demais, o ponto pode ultrapassar o
\textbf{joelho} da curva de desmagnetização --- e aí ocorre perda
\textbf{irreversível} de magnetização.
\end{itemize}
\textbf{Consequência prática decisiva:} desmontar um circuito magnético com ímã
de alnico pode danificá-lo permanentemente, pois o ponto de operação despenca.
Ímãs de neodímio, cuja curva de desmagnetização é praticamente reta até o
segundo quadrante inteiro, toleram isso muito melhor --- e é uma das razões de
sua popularidade, além do produto de energia.}
\end{questao}
